\subsection{Running the estimation}
\label{sec:pfit-run}

After the JAX program has been generated, the user starts parameter estimation
with

\begin{verbatim}
/pfit-run <session-name>
\end{verbatim}

By default, this command performs the full fitting procedure: a global search
followed by local gradient-based refinement. A refinement-only run can be
requested when a previous run has already found a suitable design point:

\begin{verbatim}
/pfit-run <session-name> gradient-only
\end{verbatim}

Before launching the optimizer, the agent performs a run pre-flight check. This
check verifies that the required files exist, that the generated JAX program is
current with respect to the checked YAML and model files, that result writing is
enabled, and, for a gradient-only run, that the seed design point is available.
If a blocking check fails, fitting is not started. The user is directed back to
the appropriate earlier step, usually \texttt{/pfit-check} or
\texttt{/pfit-jax}.

Each run is written to a new output directory, so earlier results are preserved.
At launch, the agent reports the exact run directory and a local live-view URL.
The live view shows the active optimization stage, current best loss, elapsed
time, and current best parameter values in physical units. It is intended for
monitoring progress; it is not a convergence assessment.

During fitting, the user should expect the first evaluation to include JAX
compilation, which can be slower than later iterations. For stiff systems or
large searches, the run may take substantially longer. If progress appears
abnormally slow, the agent reports the evidence from the active run and may
propose an intervention, such as stopping the current run, preserving its
artifacts, changing a numerical setting, regenerating the JAX program, and
starting a new run. Such an intervention is a restart, not a checkpoint resume,
and is applied only after the user approves the specific change.

When the run finishes, the agent reads the results from the output files rather
than inferring them from console output. The completion report includes the run
identifier, final objective value, fitted parameter values, and the locations of
the saved result tables. The agent also generates and inspects
measured-versus-fitted plots for every fitted observable and experiment. These
plots are required because the numerical objective alone does not verify that
the data columns, observables, and trajectories have been interpreted correctly.

The user should review the reported parameters and saved fit plots before
drawing conclusions. Formal assessment of convergence, residual structure,
model misspecification, and practical identifiability is deferred to the next
workflow step:

\begin{verbatim}
/pfit-diagnose <session-name> <run-id>
\end{verbatim}
